%% ---------------------------------------------------------------------------
\section{Introduction}\label{sec:intro}

Many tasks in structural engineering require the same structural system to be analysed many times: preliminary sizing and design-space exploration, code-based checks over large numbers of load combinations and variants, portfolio-scale assessment of existing buildings, optimisation and reliability analysis. Each analysis is cheap in isolation, but the number of repetitions makes finite element (FE) analysis the dominant cost of these loops. Surrogate models, which replace the FE solver by a fast approximation learned from a limited number of FE runs, have therefore a long history in engineering design~\cite{forrester2008surrogate} and reliability analysis~\cite{bichon2008egra,echard2011akmcs}, and machine learning (ML) has substantially widened the range of quantities that can be emulated~\cite{xie2020eqml}.

For structural assessment, however, a surrogate is only useful if its output can be trusted in the way an FE result is trusted. Three obstacles stand in the way. First, purely data-driven models carry no guarantee that their outputs are mechanically admissible: predicted member forces need not be in equilibrium with the applied loads, the predicted response need not scale linearly with the loads in a linear analysis, and boundary conditions are only approximately respected. An engineer who checks the reactions against the loads will immediately notice such defects. Second, accuracy reported on test data drawn from the training distribution says little about frames whose geometry or loading differ from those seen in training, which is the normal situation in practice. Third, the cost of generating training data is often treated as fixed, although the FE budget is usually the scarcest resource, and the surrogate itself does not indicate when its prediction should not be used.

Physics-informed machine learning~\cite{karniadakis2021piml} addresses the first obstacle by embedding physical knowledge in the learning model. Physics-informed neural networks penalise the residual of the governing equations in the loss~\cite{raissi2019pinn} and have been applied to solid mechanics~\cite{haghighat2021solid}; in structural dynamics, physics-guided convolutional and recurrent networks have been used to model the seismic response of buildings~\cite{zhang2020phycnn,zhang2020phylstm}. A penalty term, however, only encourages the constraint and leaves a residual violation whose size depends on the training. Constraints can instead be imposed exactly, either through the network output layer~\cite{beucler2021constraints} or through hard-constraint formulations~\cite{lu2021hard}. Graph neural networks (GNNs)~\cite{battaglia2018relational,gilmer2017mpnn} are a natural representation for structures whose topology varies: they have been used to simulate particle and mesh-based physics~\cite{sanchez2020learning,pfaff2021meshgraphnets}, to predict stress and deformation fields~\cite{maurizi2022gnn} and to build reusable surrogates for truss structures~\cite{whalen2022truss}. On the side of uncertainty, deep ensembles~\cite{lakshminarayanan2017ensembles} provide a simple estimate of epistemic uncertainty, although its quality degrades under distribution shift~\cite{ovadia2019shift}; and adaptive sampling driven by surrogate uncertainty is well established for Kriging models in structural reliability~\cite{bichon2008egra,echard2011akmcs} and more broadly as active learning~\cite{settles2009active}.

What is missing is a workflow for skeletal structures that combines these elements so that (a) the predicted member forces are statically admissible by construction rather than on average, (b) the known linear structure of the problem is exploited instead of being learned, (c) the uncertainty estimate is used to decide both where to spend FE analyses during training and which predictions to hand back to the FE solver at deployment, and (d) the result is tested against conventional regression and ML baselines on geometries and load cases that the models have not seen. This paper develops and evaluates such a workflow for planar moment-resisting frames, using the open-source PyNite FE library~\cite{pynite} for every analysis. The research question is whether a carefully validated surrogate can predict displacements and member forces accurately enough to reduce repeated FE analysis.

The contributions are:
\begin{enumerate}
\item A mechanics-constrained graph surrogate in which load-case superposition, homogeneity in the load intensities and Young's modulus, and the boundary conditions hold exactly by construction, and in which exact statics quantities (storey shears, overturning moments, column gravity loads) and classical approximations (shear-building drift, beam-to-column stiffness ratio) are supplied as inputs.
\item An equilibrium projection layer that maps any predicted member-force field to the nearest statically admissible field, with an error-reduction guarantee in its own metric (Proposition~\ref{prop:projection}) and a variance-weighted metric derived from the ensemble. The layer is model-agnostic and is also applied to the baselines.
\item An uncertainty-guided workflow that uses the ensemble spread to select the FE analyses used for training and to refer unreliable predictions to the FE solver at deployment, with the cost--accuracy trade-off quantified by risk--coverage curves.
\item A benchmark of 7{,}200 PyNite frame analyses with in-distribution and two out-of-distribution test sets (unseen geometries; unseen load patterns), on which the workflow is compared with linear and quadratic ridge regression, random forests, gradient boosting and a multilayer perceptron. The code and data generator are released with the PyNite fork.
\end{enumerate}

The remainder of the paper is organised as follows. Section~\ref{sec:data} describes the frame family and the PyNite data generation, Section~\ref{sec:method} the surrogate and the constraints, Section~\ref{sec:workflow} the uncertainty-guided workflow and Section~\ref{sec:protocol} the baselines and metrics. Results are presented in Section~\ref{sec:results} and discussed, with the limitations of the study, in Section~\ref{sec:discussion}. Section~\ref{sec:conclusions} concludes.
